\documentclass[11pt,a4paper]{article}

\usepackage{amsmath,amssymb,amsthm}
\usepackage{mathtools}
\usepackage{geometry}
\usepackage{hyperref}

\geometry{margin=2.5cm}

\newtheorem{definition}{Definition}
\newtheorem{theorem}{Theorem}
\newtheorem{lemma}{Lemma}
\newtheorem{proposition}{Proposition}
\newtheorem{corollary}{Corollary}
\newtheorem{conjecture}{Conjecture}
\newtheorem{assumption}{Assumption}
\newtheorem{remark}{Remark}

\title{Mathematical Problems in LIFA-Fuzz:\\Formulations Awaiting Proof}
\author{}
\date{}

\begin{document}
\maketitle

\section{Introduction}

This document formulates four mathematical problems arising from the
LIFA-Fuzz protocol fuzzer. Each problem is stated in self-contained
mathematical language with definitions, assumptions, and the specific
claim that requires proof. Background context is kept minimal.

\paragraph{Notation.}
$\mathbb{N} = \{0,1,2,\dots\}$, $\mathbb{E}$ denotes expectation,
$\operatorname{Var}$ denotes variance, $\lfloor\cdot\rfloor$ is the
integer floor.

\section{Boundedness and Monotonicity of the EWMA Adaptive Sampler}

The EWMA controller is a simple adaptive filter, not a feedback regulator
requiring stability analysis. Its properties are elementary
and follow directly from the formula.

\subsection{Dynamical Model}

Let $t \in \mathbb{N}$ index discrete epochs. The smoothed coverage
intensity is
\[
x_{t+1} = \delta \cdot \Delta_{t+1} + (1-\delta) \cdot x_t,
\qquad x_0 = 0,
\]
where $\delta \in (0,1)$ (fixed at $0.1$) and $\Delta_t \ge 0$ is the
coverage increment at epoch $t$. The sampling interval is
\[
k(x) = \max\!\Bigl(k_{\min},\;\Bigl\lfloor \frac{K_{\max}}{1+\theta x}\Bigr\rfloor\Bigr).
\]

\subsection{Elementary Properties}

\begin{theorem}[Boundedness]
$0 \le x_t \le \max_t \Delta_t$ for all $t$.
\end{theorem}

\begin{proof}
$x_0 = 0$. Each update is a convex combination:
$\delta \Delta_t + (1-\delta) x_{t-1}$. Since $\Delta_t \ge 0$ and
$\delta \in (0,1)$, the result follows by induction.
\end{proof}

\begin{theorem}[Monotonic response]
$\frac{\partial k}{\partial x} = \frac{-\theta K_{\max}}{(1+\theta x)^2} < 0$,
so $k$ is strictly decreasing in $x$.
\end{theorem}

\begin{proof}
Direct differentiation.
\end{proof}

The controller is thus free of oscillation or chattering by construction,
unlike step-function AIMD controllers that jump between $k=1$ and
$k = K_{\max}$.

\section{Correctness and Expected Time of the IFPS Sampler}

\subsection{Model}

Let $\mathcal{C} = \{s_1, \dots, s_N\}$ be a set of seeds.  Each seed
$s_i$ has a nonnegative integer frequency $f_i$ and an energy

\[
E(i) = \frac{1}{f_i + 1},
\qquad
E'(i) =
\begin{cases}
5 \cdot E(i) & \text{if seed $i$ is marked STATE\_NOVELTY}, \\
E(i)         & \text{otherwise}.
\end{cases}
\]

\subsection{Algorithm}

\begin{definition}[IFPS trial]
A single trial consists of:
\begin{enumerate}
  \item Pick $i \sim \operatorname{Uniform}\{1,\dots,N\}$.
  \item Accept $i$ with probability $a_i = \min(E'(i), 1)$.
\end{enumerate}
The algorithm repeats trials for at most $N+10$ attempts. If all
fail, it returns a uniformly random seed (the fallback).
\end{definition}

\subsection{Theorems}

\begin{theorem}[Correctness of acceptance-rejection]
\label{thm:ifps-correct}
Assume the fallback is never reached (its probability tends to $0$ as
$N \to \infty$).  Then the probability that seed $i$ is returned in a
given trial is

\[
P(\text{select } i) = \frac{E'(i)}{\sum_{j=1}^N E'(j)}.
\]
\end{theorem}

\begin{proof}
Elementary.  Pick $i$ uniformly: $P(\text{pick } i) = 1/N$.  Accept with
probability $a_i$.  By the acceptance-rejection lemma, conditioned on
acceptance occurring,

\[
P(\text{return } i) = \frac{(1/N) \cdot a_i}{\sum_j (1/N) \cdot a_j}
= \frac{a_i}{\sum_j a_j}.
\]

For seeds with $E'(i) \le 1$ (all seeds with $f_i \ge 4$ for novel
seeds, or any $f_i$ for normal seeds), $a_i = E'(i)$.  For novel seeds
with $f_i < 4$, $E'(i) > 1$ so $a_i = 1$, capping the acceptance
probability at unity.  In either case $a_i \propto E'(i)$ up to the cap.
\end{proof}

\begin{theorem}[Expected time]
\label{thm:ifps-time}
Let $a_{\min} = \min_i a_i$.  Since $E'(i) > 0$ for all finite $f_i$,
$a_{\min} > 0$.  The expected number of trials per successful selection
is at most $1 / a_{\min}$, a constant independent of $N$.  Moreover, the
frequency cap $f_i \le 999$ (enforced every $10^4$ sends) guarantees
$a_{\min} \ge 1/1000$.
\end{theorem}

\begin{proof}
Each trial succeeds with probability $p = \frac{1}{N} \sum_i a_i \ge a_{\min}$.
The number of trials $T$ is geometric with success probability $p$,
so $\mathbb{E}[T] = 1/p \le 1/a_{\min}$.  Under the cap,
$a_i \ge \min(1/1000, 1) = 1/1000$.
\end{proof}

\begin{remark}
Both theorems are proved; they are included for completeness.
\end{remark}

\subsection{Open Problem}

\begin{enumerate}
  \item \textbf{Non-asymptotic fallback probability.}
        Theorem~\ref{thm:ifps-correct} assumes the fallback is never
        reached.  Can we bound the probability of reaching the fallback
        after $N+10$ trials as a function of $N$ and $a_{\min}$?
        A sharper bound than $(1 - a_{\min})^{N+10}$ would be useful
        for selecting the retry limit.
\end{enumerate}

\section{Collision Bounds for Two-Level Crash Deduplication}

\subsection{Primary Signature (Exact Match)}

Let $p$ be a crash payload (a byte string).  Define

\[
\sigma_1(p) =
\begin{cases}
\operatorname{SHA256}(\texttt{"\_\_empty\_crash\_\_"} \;\|\;
    \texttt{monotonic\_ns})[0:16] & \text{if } p = \varepsilon, \\[4pt]
\operatorname{SHA256}(p)[0:16]    & \text{otherwise}.
\end{cases}
\]

The output is $16$ bytes ($128$ bits).  For distinct nonempty payloads
$p \neq p'$, the probability of a collision under the random-oracle model
for SHA256 is $2^{-128}$ per pair.

\begin{proposition}[Birthday bound]
\label{prop:collision}
After collecting $n$ distinct crash payloads, the probability of at
least one collision in $\sigma_1$ is approximately

\[
P(\text{collision}) \approx 1 - \exp\!\left(-\frac{n^2}{2^{129}}\right).
\]

For $n \le 2^{32}$, this gives $P(\text{collision}) \le 2^{-65}$.
\end{proposition}

\begin{proof}
Classical birthday problem with $N = 2^{128}$ bins.
\end{proof}

\subsection{Structural Signature (Similarity Grouping)}

Define the structural fingerprint

\[
\sigma_2(p) = \operatorname{XXH64}\bigl(
    \operatorname{head}_{24}(p) \;\|\;
    \operatorname{len}_{32\text{be}}(p) \;\|\;
    \operatorname{fold}(p)
\bigr)[0:6],
\]

where $\operatorname{head}_{24}(p) = p[0:24]$,
$\operatorname{len}_{32\text{be}}(p)$ is the 4-byte big-endian length,
and $\operatorname{fold}(p)$ is a content summary:

\[
\operatorname{fold}(p) =
\begin{cases}
p & |p| < 32, \\[4pt]
\operatorname{mid}_8(p) \;\|\;
    \displaystyle\bigoplus_{c \in \operatorname{chunks}_8(p)} c
    & |p| \ge 32,
\end{cases}
\]

with $\operatorname{mid}_8(p) = p[|p|/2 : |p|/2+8]$ and XOR fold of
8-byte big-endian chunks (final chunk zero-padded).  The output is
$6$ bytes ($48$ bits).

\begin{proposition}
\label{prop:structural-collision}
Under the random-oracle model for XXH64, the pairwise collision
probability for $\sigma_2$ is $2^{-48}$, and after $n$ payloads,

\[
P(\text{collision}) \approx 1 - \exp\!\left(-\frac{n^2}{2^{49}}\right).
\]
\end{proposition}

\subsection{Open Problem}

\begin{enumerate}
  \item \textbf{Structural vs.\ random collisions.}
        Proposition~\ref{prop:structural-collision} treats XXH64 as a
        random oracle.  In practice, two structurally similar payloads
        (differing by one byte) may produce correlated fingerprints.
        Can we bound the probability that $\sigma_2$ collides for
        \emph{structurally distinct} payloads, accounting for the
        specific construction of $\operatorname{fold}$?
\end{enumerate}

\section{Probabilistic Protocol State Machine Inference}

\subsection{Formal Model of a Protocol}

\begin{definition}[Protocol]
A protocol is a tuple $(\mathcal{S}, \mathcal{A}, \delta, s_0)$ where:
\begin{itemize}
  \item $\mathcal{S}$ is a finite set of states.
  \item $\mathcal{A}$ is a finite set of message types (the alphabet).
  \item $\delta: \mathcal{S} \times \mathcal{A} \to \mathcal{S}$ is the
        (deterministic) transition function.
  \item $s_0 \in \mathcal{S}$ is the initial state.
\end{itemize}
\end{definition}

Each captured packet $p$ is generated by some state $s$ and carries an
(unknown) message type $a \in \mathcal{A}$.  The observed byte string
$p$ is a noisy rendering of $a$; the precise mapping
$\nu: \mathcal{A} \to \{0,1\}^*$ (from message type to wire format) is
unknown.

\subsection{Inference Pipeline}

The pipeline proceeds as follows:

\begin{enumerate}
  \item \textbf{Message-unit extraction.}
        Extract all length-3 substrings (``units'') from the first 12
        bytes of each packet.  Split the corpus in half and retain only
        units appearing at least $\lambda$ times in both halves, where
        $\lambda$ is the smallest threshold such that the two empirical
        frequency distributions pass a two-sample Kolmogorov--Smirnov
        test at significance $1 - \alpha < 10^{-8}$.
  \item \textbf{Format message reconstruction.}
        For each packet, greedily extend from offset 0 using candidate
        units from step~1.  The maximal prefix composed entirely of
        candidate units is the packet's \emph{format message} $m$.
  \item \textbf{Clustering.}
        Represent each format message $m$ by its byte-frequency histogram
        $\mathbf{v}(m) \in \mathbb{N}^{256}$.  Cluster using
        Partitioning Around Medoids (PAM) under Jaccard distance on the
        support sets $\{b : v_b > 0\}$.  Select $k$ maximising the
        Dunn index.
  \item \textbf{State machine.}
        Each cluster medoid becomes a state.  Transition probabilities
        are estimated from observed sequences.
\end{enumerate}

\subsection{Open Problem}

The central question is: under what conditions on the protocol
$(\mathcal{S}, \mathcal{A}, \delta, s_0)$ and the encoding $\nu$ does
this pipeline recover the true state machine?

\begin{conjecture}[Sufficient condition for recovery]
\label{conj:ppsm}
Assume:
\begin{enumerate}
  \item \textbf{Distinct keyword prefixes.}
        For each message type $a \in \mathcal{A}$, the first 12 bytes
        of $\nu(a)$ contain at least one length-3 substring not present
        in the first 12 bytes of $\nu(a')$ for any $a' \neq a$.
        (Equivalently, the set of 3-grams in the first 12 bytes
        uniquely identifies each message type.)
  \item \textbf{Frequency separation.}
        There exists a threshold $\lambda$ such that protocol message
        types appear in every session while random byte sequences
        (payload data) do not, and the K-S test at $\alpha < 10^{-8}$
        distinguishes the two classes.
  \item \textbf{Homogeneous byte distributions.}
        Packets of the same message type $a$ have byte-frequency
        histograms $\mathbf{v}(\nu(a))$ whose Jaccard distance to each
        other is strictly less than their distance to histograms of
        any other message type $a' \neq a$.
\end{enumerate}
Then the pipeline recovers a state graph isomorphic to the true protocol
state machine, up to relabelling of states.
\end{conjecture}

\begin{remark}
Assumptions 1--3 are strong.  In practice some message types may share
prefixes (e.g., ``220 `` and ``230 `` both begin with ``2`` and ``0``).
The goal is to determine which assumptions are necessary and which can
be relaxed, and to bound the recovery error when they are violated.
\end{remark}

\section{Summary of Open Problems}

\begin{enumerate}
  \item Sharper drift bounds and convergence rate for the EWMA closed-loop
        system (Section~\ref{conj:ultimate} and its open problems).
  \item Non-asymptotic fallback probability for the IFPS sampler
        (Section~\ref{thm:ifps-correct}).
  \item Structural collision bound for the XXH64-based crash fingerprint
        (Section~\ref{prop:structural-collision}).
  \item Sufficient conditions for protocol state machine recovery
        (Conjecture~\ref{conj:ppsm}) and relaxation thereof.
\end{enumerate}

\end{document}
